\documentclass[12pt,twoside,book]{article}
\usepackage{docmute}

\input{settings}

\begin{document}

\begin{titlepage}

\begin{center}

\begin{titlelineskip}

\vskip .75in

{\Large \it

 Ph.D. Thesis

}

\vskip 1.0in

{\Huge \bf

  Probing Electroweakly Interacting\\
  Massive Particles\\
  with Drell-Yan Process\\
  at 100 TeV Colliders

}

\vskip .75in

{\Large
So Chigusa
}\\

{\Large \it Department of Physics}

\vfill

\centering
\includegraphics[width=0.6\hsize]{univ.eps}

\vskip 0.35in

{\Large December 2019}

\end{titlelineskip}

\end{center}

\vspace{0.5in}

\end{titlepage}

\thispagestyle{empty}

\vspace*{1.50in}

\begin{abstract}
 \rem{To be written}
\end{abstract}

\clearpage

\thispagestyle{empty}

\vspace*{.75in}

{\LARGE \bf
Acknowledgments
}

\vskip .50in

\rem{To be written}

% I would like to thank my supervisor, Takeo Moroi, who provided
% stimulating discussions, helpful comments, fruitful suggestions and
% collaboration.  I am also grateful to Koichi Hamaguchi, Kazunori
% Nakayama, Natsumi Nagata and all the members of the particle physics
% group at the University of Tokyo for their hospitality.  My work is
% supported by the Program for Leading Graduate Schools, MEXT, Japan.
% Finally, I would like to express my gratitude to my family, friends and
% all things which allow me to enjoy studying physics.

\clearpage

\thispagestyle{empty}

\tableofcontents

\clearpage

\renewcommand{\thepage}{\arabic{page}}
\setcounter{page}{1}

%%%%%%%%%%%%%%%%%%%%%%%%%%%%%%%%%%%%%%%%%%%%%%%%%%%%%%%%%%%%%%%%%%%%%%%%%%%%%%%
\section[Introduction]{Introduction}
\setcounter{equation}{0}
%%%%%%%%%%%%%%%%%%%%%%%%%%%%%%%%%%%%%%%%%%%%%%%%%%%%%%%%%%%%%%%%%%%%%%%%%%%%%%%

\vskip 0.1in

\rem{Unit: $\hbar=c=k_B=1$}

\rem{Convention: dot as time derivative}

\rem{Definition of ``SM''}

\rem{Definition of ``WIMP''}

% Erase biblio in input file before turning on
% \input{WIMP/wimp.tex}

\clearpage

\begin{appendices}

%%%%%%%%%%%%%%%%%%%%%%%%%%%%%%%%%%%%%%%%%%%%%%%%%%%%%%%%%%%%%%%%%%%%%%%%%%%%%%%
\section{Review of supersymmetry}
%%%%%%%%%%%%%%%%%%%%%%%%%%%%%%%%%%%%%%%%%%%%%%%%%%%%%%%%%%%%%%%%%%%%%%%%%%%%%%%

In this appendix, we briefly review the $\mathcal{N}=1$ supersymmetry,
which is an essential element of the MSSM reviewed in
Sec.~\ref{sec:mssm}.  Our argument is based on~\cite{Wess:320631,
Martin:1997ns}.

The $\mathcal{N}=1$ supersymmetry is

First example is the MSSM, extension of the SM with the so-called
$\mathcal{N}=1$ supersymmetry (SUSY)~\cite{Wess:320631, Martin:1997ns}
that relates a bosonic particle and a fermionic particle.  The
supersymmetry transformations for a complex scalar $\phi$ and its
``superpartner'' Weyl fermion $\psi$ are defined as
\begin{align}
 \delta \phi = \left( \epsilon \psi \right),
 ~&~
 \delta \phi^{*} = \left( \epsilon^\dagger \psi^\dagger \right),\\
 \delta \psi = -i \left(\sigma^\mu \epsilon^\dagger \right) \partial_\mu \phi,
 ~&~
 \delta \psi^\dagger = i \left(\epsilon \sigma^\mu \right) \partial_\mu \phi^{*},
\end{align}
where $\sigma^\mu \equiv (\bm{1}, \bm{\sigma})$ with $\bm{\sigma}$
being Pauli matrices, while $\epsilon$ is an anti-commuting Weyl
fermionic object that parameterizes the SUSY transformation.  The
summation over the spinor indices is assumed inside each parenthesis.
These transformations, if denoted by operators $\epsilon Q$ and
$\epsilon^\dagger Q^\dagger$, are known to form a closed algebra
\begin{align}
 \left[ Q, Q^\dagger \right] &= 2 i \sigma^\mu \partial_\mu,\\
 \left[ Q, Q \right] &= \left[ Q^\dagger, Q^\dagger \right] = 0,
\end{align}
when fields are on-shell.\footnote{
%%
In order for the algebra to be closed off-shell, one can introduce a
new scalar field $F$ without a kinetic term that is often called as an
\textit{auxiliary} field.  $F$ works as a Lagrange multiplier whose
equation of motion }

\end{appendices}

\clearpage

\bibliographystyle{elsarticle-num}
\bibliography{phd}

\end{document}
